\documentclass[addpoints]{exam}
\usepackage[utf8]{inputenc}

\usepackage[12pt]{extsizes}
\usepackage[margin=0.8in]{geometry}

\usepackage{times}
\usepackage{color}
\usepackage{graphicx}

\title{CSE114A, Spring 2025: Midterm Exam 1}
\author{Instructor: Lindsey Kuper}
\date{May 9, 2025}

\renewcommand{\thechoice}{\alph{choice}}
\renewcommand{\choicelabel}{(\thechoice)}

\renewcommand{\thepartno}{\alph{partno}}
\renewcommand{\partlabel}{\thepartno.}

\setlength\fillinlinelength{2in}
\setlength\answerclearance{0.3ex}

\CorrectChoiceEmphasis{\color{red}}
\shadedsolutions

\begin{document}

\maketitle

\begin{coverpages}

\bigskip  

\noindent Student name: \hrulefill

\bigskip
\bigskip

\noindent CruzID: \hrulefill @ucsc.edu

\bigskip
\bigskip

\noindent This exam has \numquestions~questions and \numpoints~total points.

\bigskip
\bigskip

\noindent \textbf{Instructions}

\begin{itemize}
\item Please write directly on the exam.
\item For short answer questions, please write your answer in the provided boxes.  You can use space outside of the boxes as scratch space, but we won't see or grade it.
\item For multiple choice questions, please completely fill in the circle for the correct choice.
\item \textbf{You have 65 minutes to complete this exam.} You may leave when you are finished.
\item This exam is \textbf{closed book}.  You may use one double-sided page of notes, but no other materials.
\item Avoid seeing anyone else's work or allowing yours to be seen.
\item Please, no talking.  No notes, books, laptops, phones, or other electronic devices. Do not communicate with anyone but an exam proctor.
\item To ensure fairness (and the appearance thereof),
  \textbf{proctors will not answer questions about the content of the exam}. If you are unsure 
  of how to interpret a problem description, state your interpretation clearly and concisely. 
  \textit{Reasonable interpretations} will be taken into account by graders.
\item \textbf{We will give partial credit for partially correct answers when it makes sense to do so.} A partially correct answer is better than leaving an answer blank.
\end{itemize}
  
\bigskip

\noindent Good luck!

\newpage

This page is for your use as scratch space. Anything you write here will be ungraded.

\newpage  

\end{coverpages}

\subsection*{\textbf{Part 1: Lambda Calculus (68 points)}}

\begin{questions}

\question[6] A lambda calculus expression is in \emph{normal form} if it cannot be further reduced.  Evaluate the following lambda calculus expression to normal form using a series of $\beta$-reduction and/or $\alpha$-renaming steps.  Start each line with \texttt{=b>} or \texttt{=a>}, as if you were using Elsa, and do just one $\beta$-reduction step or $\alpha$-renaming step per line.

There may be multiple correct ways to reduce the given expression.  A correct solution is any solution that Elsa would accept as correct.

\begin{verbatim}
\y -> (\x z -> z x) (\r q -> r q) (\s q -> s q) y
\end{verbatim}
\begin{solutionorbox}[\stretch{1}]
Here is a possible answer (with comments for explanation).
\begin{verbatim}
  \y -> (\x z -> z x) (\r q -> r q) (\s q -> s q) y
  -- Reduce the redex `(\x z -> z x) (\r q -> r q)`
  -- (this is the only redex at this step)
  =b> \y -> (\z -> z (\r q -> r q)) (\s q -> s q) y
  -- Reduce the redex `(\z -> z (\r q -> r q)) (\s q -> s q)`
  -- (this is the only redex at this step)
  =b> \y -> (\s q -> s q) (\r q -> r q) y
  -- Reduce the redex `(\s q -> s q) (\r q -> r q)`
  -- (this is the only redex at this step)
  =b> \y -> (\q -> (\r q -> r q) q) y
  -- Rename `q` to `p` so we can take a beta-step on `(\r q -> r q) q`
  =a> \y -> (\q -> (\r p -> r p) q) y
  -- Reduce the redex `(\r p -> r p) q`
  -- (this is one of two possible redexes here)
  =b> \y -> (\q -> (\p -> q p)) y
  -- Reduce the redex `(\q -> (\p -> q p)) y`
  -- (this is the only redex at this step)
  =b> \y -> (\p -> y p)
  -- We now have a value; no more beta-steps are possible
\end{verbatim}
In the above renaming (\verb|=a>|) step, any unused variable is fine; it doesn’t have to be \verb|p|.

Alternatively, once we have \verb|\y -> (\q -> (\r q -> r q) q) y|, we could end the reduction like this, which wouldn't require renaming:
\begin{verbatim}
  =b> \y -> (\q -> (\r q -> r q) q) y
  -- Reduce the redex `(\q -> (\r q -> r q) q) y`
  -- (this is one of two possible redexes here)
  =b> \y -> (\r q -> r q) y
  -- Reduce the redex `(\r q -> r q) y`
  -- (this is the only redex at this step)
  =b> \y -> (\q -> y q)
  -- We now have a value; no more beta-steps are possible
\end{verbatim}

(Notice that the expressions \verb|\y -> (\p -> y p)| and \verb|\y -> (\q -> y q)| are alpha-equivalent.)
\end{solutionorbox}

\newpage

Refer to the lambda calculus reference at the end of the exam for definitions used in this section.

\question[6] Which of the following lambda calculus expressions is in \textbf{normal form}? (Choose only one answer.)

\begin{checkboxes}
  \choice \verb|(\x -> x) y|
  \choice \verb|\z -> (\x -> x) (\y -> y)|
  \choice \verb|\z -> (\y -> y) z|
  \CorrectChoice \verb|\z -> (\y -> y)|
  \choice \verb|(\z -> z z) (\z -> z z)|
\end{checkboxes}

\question[6] Which of the following lambda calculus expressions evaluates to \texttt{FALSE}? (Choose only one answer.)

\begin{checkboxes}
  \choice \texttt{FALSE FALSE TRUE}
  \choice \texttt{FALSE FALSE}
  \choice \texttt{ITE FALSE FALSE}
  \choice \texttt{ITE FALSE FALSE TRUE}
  \CorrectChoice \texttt{FALSE TRUE FALSE}
\end{checkboxes}

\question[6] Which of the following lambda calculus expressions evaluates to \texttt{ZERO}? (Choose only one answer.)

\begin{checkboxes}
  \CorrectChoice \texttt{SND (PAIR (PAIR ONE ZERO) ZERO)}
  \choice \texttt{TRUE ZERO}
  \choice \texttt{TRUE ONE ZERO}
  \choice \texttt{FALSE ZERO ONE}
  \choice \texttt{FALSE ZERO}
\end{checkboxes}

\question[6] What is the set of variables that \textbf{occur free} in the following lambda calculus expression? (Choose only one answer.)

\begin{verbatim}
(\x y -> y) z (\y -> x z y)
\end{verbatim}

\begin{checkboxes}
  \choice No variables occur free
  \choice \texttt{z}
  \choice \texttt{x, y}
  \CorrectChoice \texttt{x, z}
  \choice \texttt{x, y, z}
\end{checkboxes}

\newpage

For the rest of this section, you may use any of the helper functions defined in the provided lambda calculus reference at the end of the exam.

\question[7] Define a lambda calculus function \texttt{PAIRSUMINCR} whose input is a \texttt{PAIR} of Church numerals (or expressions that evaluate to Church numerals) and whose output is the Church numeral with value one greater than the sum of the two input numerals.  For example, in Elsa:

\begin{verbatim}
PAIRSUMINCR (PAIR ZERO ZERO) =~> ONE
PAIRSUMINCR (PAIR ONE ZERO) =~> TWO
PAIRSUMINCR (PAIR TWO ONE) =~> FOUR
PAIRSUMINCR (PAIR ONE (DECR (DECR TWO))) =~> TWO
\end{verbatim}

You may assume that \texttt{PAIRSUMINCR} is only called with \texttt{PAIR}s of expressions that evaluate to Church numerals.

\begin{verbatim}
let PAIRSUMINCR = 
\end{verbatim}
\begin{solutionorbox}[2in]
One simple correct answer is:
\begin{verbatim}
\p -> INCR (ADD (FST p) (SND p))
\end{verbatim}
Other correct answers are also possible.
\end{solutionorbox}

\newpage

The next two questions will involve a lambda calculus encoding of \textbf{lists}.  We can encode lists in lambda calculus using pairs, as follows:

\begin{verbatim}
let NIL   = \x -> TRUE
let CONS  = PAIR
let HEAD  = FST
let TAIL  = SND
let ISNIL = \lst -> lst (\h t -> FALSE)
\end{verbatim}

The list constructors are \texttt{NIL} and \texttt{CONS}.  \texttt{NIL} is the empty list, and a non-empty list consists of a pair of a list element and a list.

For instance, \texttt{CONS ONE (CONS TWO (CONS THREE NIL))} is the lambda calculus equivalent of the Haskell list \texttt{[1, 2, 3]},\\
which is just syntactic sugar for \texttt{1 : (2 : (3 : []))}.

The \texttt{ISNIL} function takes a list and returns \texttt{TRUE} if the list is \texttt{NIL} and \texttt{FALSE} otherwise.

You can (and should!) use the \texttt{NIL}, \texttt{CONS}, \texttt{HEAD}, \texttt{TAIL}, and \texttt{ISNIL} functions in your answers to the next two questions.

\question[7]

Define a lambda calculus function \texttt{INCRSECOND} that takes a list of Church numerals, and returns the second element of the list, incremented by one.  You may assume that the argument to \texttt{INCRSECOND} is a list of \emph{at least} two Church numerals.  Any list elements other than the second one should be ignored.  For example, in Elsa:

\begin{verbatim}
  INCRSECOND (CONS TWO (CONS TWO NIL)) =~> THREE
  INCRSECOND (CONS ZERO (CONS TWO NIL)) =~> THREE
  INCRSECOND (CONS ZERO (CONS ONE (CONS TWO NIL))) =~> TWO
  INCRSECOND (CONS ZERO (CONS ZERO (CONS THREE NIL))) =~> ONE
  INCRSECOND (CONS THREE (CONS ONE (CONS ZERO NIL))) =~> TWO
\end{verbatim}

\begin{verbatim}
let INCRSECOND = 
\end{verbatim}
\begin{solutionorbox}[2in]
One simple correct answer is:
\begin{verbatim}
\lst -> INCR (HEAD (TAIL lst))
\end{verbatim}
\end{solutionorbox}

\newpage

\question Define a lambda calculus function \texttt{DOUBLEALL} that takes a list of Church numerals, and returns a list with the same elements except that each element has been \emph{doubled}.  For example:

\begin{verbatim}
  DOUBLEALL NIL =*> NIL

  DOUBLEALL (CONS ONE NIL) =*> CONS TWO NIL

  DOUBLEALL (CONS ONE (CONS TWO NIL)) =*> CONS TWO (CONS FOUR NIL)

  DOUBLEALL (CONS TWO (CONS ZERO (CONS ONE NIL))) =*>
    CONS FOUR (CONS ZERO (CONS TWO NIL))
\end{verbatim}

You may assume that \texttt{DOUBLEALL} is called with a list constructed with \texttt{NIL} or \texttt{CONS}, and that all elements in the list are Church numerals.  You must use recursion for full credit.  Hint: You can double a list element by \texttt{ADD}ing it to itself.

\begin{verbatim}
let DOUBLEALL1 = \rec -> \lst -> ITE _____(part 8(a))_______
                                     _____(part 8(b))_______
                                     _____(part 8(c))_______
                      
let DOUBLEALL = ______(part 8(d))_______
\end{verbatim}

\begin{parts} 
  \part[6] 8(a):
  \begin{solutionorbox}[\stretch{4}]
    This is where we check a condition to know whether we are in the base case or not.
    \texttt{ISNIL lst} is a correct answer.
  \end{solutionorbox}

  \part[6] 8(b):
  \begin{solutionorbox}[\stretch{4}]
This is the base case.  \texttt{NIL} is a correct answer.
  \end{solutionorbox}

  \part[6] 8(c):
  \begin{solutionorbox}[\stretch{4}]
This is the recursive case.  \texttt{CONS (ADD (HEAD lst) (HEAD lst)) (rec (TAIL lst))} is a correct answer.
  \end{solutionorbox}

  \part[6] 8(d):
  \begin{solutionorbox}[\stretch{4}]
\begin{verbatim}
Y DOUBLEALL1
\end{verbatim}
  \end{solutionorbox}
\end{parts}

\newpage

\subsection*{\textbf{Part 2: Haskell (32 points)}}

The Haskell reference at the end of the exam has information about library functions used in this section.

\question[6] What is the \textbf{type} of the following Haskell expression?

\begin{verbatim}
["rainbow" ++ "sprinkles"] ++ ["mo" ++ "larry"]
\end{verbatim}

\begin{checkboxes}
  \choice This expression is ill-typed; it doesn't have a type
  \choice \verb|String|
  \CorrectChoice \verb|[String]|
  \choice \verb|[[String]]|
  \choice \verb|[String] -> [String] -> [String]|
\end{checkboxes}

\question[6] What is the \textbf{type} of the following Haskell expression?

\begin{verbatim}
(\x -> if x then ["mo", "sam"] else ["the creature"]) False
\end{verbatim}

\begin{checkboxes}
  \choice This expression is ill-typed; it doesn't have a type
  \choice \verb|Bool -> [String]|
  \CorrectChoice \verb|[String]|
  \choice \verb|Bool -> String|
  \choice \verb|String|
\end{checkboxes}

\question[6] What is the \textbf{type} of the following Haskell expression?

\begin{verbatim}
\x -> "phoebe" : [x]
\end{verbatim}

\begin{checkboxes}
  \choice This expression is ill-typed; it doesn't have a type
  \choice \verb|String -> String|
  \choice \verb|[String] -> String|
  \CorrectChoice \verb|String -> [String]|
  \choice \verb|[String] -> [String]|
\end{checkboxes}

\newpage

For the next two questions, you can use the library functions provided in the Haskell reference at the end of the exam.

Haskell's standard library provides a list data type with two constructors, \verb|[]| (pronounced ``nil'') and \verb|(:)| (pronounced ``cons'').  For instance, the list \verb|[True, False, False]| can be constructed with calls to the list constructors:

\begin{verbatim}
True : (False : (False : []))
\end{verbatim}

But we could also define our own list data type.  For instance, here is a type of lists of \verb|Bool|s:

\begin{verbatim}
data BooList = BNil | BCons Bool BooList
\end{verbatim}

To construct a \verb|BooList| with the elements \verb|True|, \verb|False|, and \verb|False|, in that order, we would use the \verb|BooList| constructors as follows:

\begin{verbatim}
BCons True (BCons False (BCons False BNil))
\end{verbatim}

\question[7] Write a Haskell function \verb|booListFlip| that takes a \verb|BooList| and returns the same list but with every element negated.  For example, in GHCi:

\begin{verbatim}
ghci> booListFlip Nil
Nil

ghci> booListFlip (BCons True (BCons False (BCons False BNil)))
BCons False (BCons True (BCons True BNil))
\end{verbatim}

The type signature is provided for you.  Hint: use pattern matching to deal with the two \verb|BooList| constructors, \verb|BNil| and \verb|BCons|, and use \verb|not| to negate a value of \verb|Bool| type.

\begin{verbatim}
booListFlip :: BooList -> BooList
\end{verbatim}
\begin{solutionorbox}[3in]
\begin{verbatim}
booListFlip BNil         = BNil
booListFlip (BCons x xs) = BCons (not x) (booListFlip xs)
\end{verbatim}
\end{solutionorbox}

\newpage

\question[7] Write a Haskell function \verb|booListAnd| that takes a \verb|BooList| and returns \verb|True| if all of the \verb|BooList|'s elements are \verb|True| and \verb|False| otherwise.  For example, in GHCi:

\begin{verbatim}
ghci> booListAnd BNil
True

ghci> booListAnd (BCons True (BCons False BNil))
False

ghci> booListAnd (BCons True (BCons True BNil))
True

ghci> booListAnd (BCons False (BCons False (BCons True BNil)))
False
\end{verbatim}

The type signature is provided for you.

Hint: use pattern matching to deal with the two \verb|BooList| constructors, and use the \verb|(&&)| operator to combine values of \verb|Bool| type.

\begin{verbatim}
booListAnd :: BooList -> Bool
\end{verbatim}
\begin{solutionorbox}[3in]
\begin{verbatim}
booListAnd :: BooList -> Bool
booListAnd BNil         = True
booListAnd (BCons x xs) = x && booListAnd xs
\end{verbatim}
\end{solutionorbox}

\end{questions}

\newpage

\subsection*{Lambda Calculus Reference}

\begin{verbatim}
-- Church numerals
let ZERO  = \f x -> x
let ONE   = \f x -> f x
let TWO   = \f x -> f (f x)
let THREE = \f x -> f (f (f x))
let FOUR  = \f x -> f (f (f (f x)))

-- Booleans
let TRUE  = \x y -> x
let FALSE = \x y -> y
let ITE   = \b x y -> b x y
let AND   = \b1 b2 -> ITE b1 b2 FALSE

-- Pairs
let PAIR   = \x y -> (\b -> ITE b x y)
let FST    = \p -> p TRUE
let SND    = \p -> p FALSE

-- Lists
let NIL   = \x -> TRUE
let CONS  = PAIR
let HEAD  = FST
let TAIL  = SND
let ISNIL = \lst -> lst (\h t -> FALSE)

-- Arithmetic
let INCR  = \n f x -> f (n f x)
let ADD   = \n m -> n INCR m

-- The definitions of DECR, SUB, ISZ, and EQL are elided
-- but you can still use them:
let DECR  = \n ->   -- (decrement n by one)
let SUB   = \n m -> -- (subtract m from n)
let ISZ   = \n ->   -- (return TRUE if n == 0 and FALSE otherwise)
let EQL   = \n m -> -- (return TRUE if n == m and FALSE otherwise)

-- Note: Since ZERO is the smallest Church numeral,
-- calls to DECR and SUB bottom out at ZERO.
-- For example, DECR ZERO evaluates to ZERO,
-- and SUB TWO THREE evaluates to ZERO.

-- The Y combinator
let Y = \step -> (\x -> step (x x)) (\x -> step (x x))
\end{verbatim}

\newpage

\subsection*{Haskell Reference}

\begin{itemize}

\item
\begin{verbatim}
(:) :: a -> [a] -> [a]
\end{verbatim}
The list constructor operator.  Takes an element and a list, and constructs a new list with the specified element as its head and the specified list as its tail.
\begin{verbatim}
> 3 : [4, 5]
[3, 4, 5]
> 3 : (4 : (5 : []))
[3, 4, 5]
\end{verbatim}

\item
\begin{verbatim}
(+) :: Num a => a -> a -> a
\end{verbatim}
Returns the sum of its two arguments, e.g.,
\begin{verbatim}
> 3 + 4
7
\end{verbatim}

\item
\begin{verbatim}
not :: Bool -> Bool
\end{verbatim}
Boolean negation.
\begin{verbatim}
> not True
False
> not (not True)
True
\end{verbatim}

\item
\begin{verbatim}
(&&) :: Bool -> Bool -> Bool
\end{verbatim}
The logical ``and'' operation, e.g.,
\begin{verbatim}
> True && False
False
> True && (not False)
True
\end{verbatim}

\item
\begin{verbatim}
(++) :: [a] -> [a] -> [a]
\end{verbatim}
Appends two lists, e.g.,
\begin{verbatim}
> [1,2,3] ++ [4,5]
[1,2,3,4,5]
> "apple" ++ "orange"
"appleorange"
\end{verbatim}

\end{itemize}

\end{document}
